% ============================================================
% Secao 8: Modulos de Extensao (9 modulos)
% ============================================================

\section{M\'odulos de Extens\~ao}
\label{sec:extensoes}

Al\'em do n\'ucleo epist\^emico (DRM, MAD, VI, MPC), o \aletheion{} inclui
9 m\'odulos de extens\~ao organizados em 3 tiers. Cada m\'odulo \'e um
\texttt{nn.Module} trein\'avel com flag \texttt{enable\_*} para
ativa\c{c}\~ao condicional.

% -------------------------------------------------------
\subsection{Tier 1: Alto Impacto, Baixo Custo (~1.7K params)}
% -------------------------------------------------------

\subsubsection{EidosDecay}
\label{sec:eidos}

O \textit{EidosDecay} monitora e corrige o balanceamento dos 5 eixos do manifold.
Eixos sobre-representados sofrem \textit{decay} (atenua\c{c}\~ao) e eixos
sub-representados recebem \textit{reinforce} (refor\c{c}o):

\begin{align}
    r_k &= \frac{\text{std}(c_k)}{\max_j \text{std}(c_j) + \epsilon} \label{eq:eidos_ratio} \\
    w_k^{\text{decay}} &= \sigma\big(\text{MLP}_{\text{decay}}(\bm{r})\big)_k \cdot \alpha_{\text{base}} \\
    w_k^{\text{reinforce}} &= \sigma\big(\text{MLP}_{\text{reinforce}}(\bm{r})\big)_k \cdot \beta_{\text{base}}
\end{align}

com $\alpha_{\text{base}} = 0.95$ e $\beta_{\text{base}} = 1.05$. O resultado
\'e um vetor \texttt{axis\_balance} $\in \R^{B \times T \times 5}$ que indica
o peso relativo de cada eixo.

\paragraph{Dream Mode.} Durante fases de ``sonho'' (treinamento especial),
a intensidade \'e amplificada por $\delta_{\text{dream}} = 3.0$:

\begin{equation}
    w_k^{\text{reinforce}} = \sigma(\cdot) \cdot \beta_{\text{base}} \cdot \delta_{\text{dream}}
\end{equation}

\paragraph{Par\^ametros:} ~368 (2 MLPs $5 \to 16 \to 5$ + 1 proje\c{c}\~ao $5 \to 1$).

\subsubsection{Filosofia3: Conflito $\phi$-$\psi$}
\label{sec:filosofia3}

O m\'odulo \textit{Filosofia3} detecta e quantifica conflitos entre
o estado epist\^emico ($\phi$) e o estado pragm\'atico ($\psi$) do modelo.

O conflito \'e medido em duas vias:

\paragraph{Via Anal\'itica.}
Computa deltas entre tokens adjacentes e calcula similaridade cosseno:

\begin{align}
    \Delta\phi_t &= \phi_{t} - \phi_{t-1} \in \R^4 \\
    \Delta\psi_t &= (\kappa_t - \kappa_{t-1}) \cdot \bm{p} \in \R^4
\end{align}

onde $\bm{p} = (0.1, 0.1, 0.1, 0.7)$ \'e a proje\c{c}\~ao de qualidade.
O conflito anal\'itico \'e:

\begin{equation}
    \gamma_{\text{anal}} = \frac{1}{2}\left(1 - \frac{\Delta\phi \cdot \Delta\psi}{\|\Delta\phi\| \|\Delta\psi\| + \epsilon}\right) \in [0, 1]
\end{equation}

\paragraph{Via Aprendida.}
Uma MLP processa a concatena\c{c}\~ao $[\Delta\phi; \Delta\psi] \in \R^8$:

\begin{equation}
    \gamma_{\text{learn}} = \sigma\big(\text{MLP}_{8 \to 16 \to 1}([\Delta\phi; \Delta\psi])\big)
\end{equation}

\paragraph{Blend Final.}
\begin{equation}
    \gamma = 0.7 \cdot \gamma_{\text{anal}} + 0.3 \cdot \gamma_{\text{learn}}
\end{equation}

\paragraph{Modos de Opera\c{c}\~ao.}
Uma segunda MLP classifica o modo em 4 categorias via softmax:

\begin{equation}
    \bm{m} = \text{softmax}\big(\text{MLP}_{8 \to 16 \to 4}([\Delta\phi; \Delta\psi])\big)
\end{equation}

\begin{enumerate}
    \item \textbf{ALIGNED} ($m_0$): $\phi$ e $\psi$ concordam
    \item \textbf{CONFLICT\_TOLERATED} ($m_1$): Conflito aceit\'avel
    \item \textbf{SIGNAL\_HUMAN} ($m_2$): Conflito requer sinaliza\c{c}\~ao
    \item \textbf{RECOVERY} ($m_3$): Conflito cr\'itico
\end{enumerate}

\paragraph{Par\^ametros:} ~373 (2 MLPs + buffers).

\subsubsection{SelfModel: Auto-Modelo Cognitivo}
\label{sec:selfmodel}

O \textit{SelfModel} mant\'em um modelo interno do estado cognitivo,
produzindo 4 sinais:

\begin{align}
    \text{mood} &= \tanh\big(\text{MLP}_{3 \to 32 \to 1}(\bar{\kappa}, \phi, p_{\text{success}})\big) \in [-1, 1] \\
    \text{curiosity} &= \sigma\big(\text{MLP}_{1 \to 8 \to 1}(q_2)\big) \in [0, 1] \\
    \text{energy} &= \sigma(\bm{w}^\top \bh) \cdot \text{decay}(t) \in [0, 1] \\
    \text{drives} &= \sigma\big(\text{MLP}_{3 \to 8 \to 3}(\text{curiosity}, \text{mood}, \text{energy})\big) \in [0, 1]^3
\end{align}

onde $\text{decay}(t) = 1 - 0.3 \cdot t / T$ simula a deple\c{c}\~ao de
energia ao longo da sequ\^encia, e os 3 drives s\~ao: curiosidade,
maestria e efici\^encia.

\paragraph{Par\^ametros:} ~982 (dependente de $d_{\text{model}}$).

% -------------------------------------------------------
\subsection{Tier 2: M\'edio Impacto, M\'edio Custo (~112K params)}
% -------------------------------------------------------

\subsubsection{Grounding: Classifica\c{c}\~ao de Tarefas e Ambiguidade}
\label{sec:grounding}

\paragraph{TaskClassificationHead.}
Classifica cada token em 9 tipos de tarefa via MLP:

\begin{equation}
    \bm{p}_{\text{task}} = \text{softmax}\big(\text{MLP}_{d \to 64 \to 9}(\bh)\big) \in \R^{B \times T \times 9}
\end{equation}

Os 9 tipos s\~ao: FACTUAL, ANALYTICAL, CREATIVE, CONVERSATIONAL,
INSTRUCTIONAL, MATHEMATICAL, CODING, SUMMARIZATION, EPISTEMIC.

\paragraph{AmbiguityHead.}
Detecta ambiguidade em dois n\'iveis:

\begin{align}
    \text{level} &= \sigma\big(\text{MLP}_{d \to 32 \to 1}(\bh)\big) \in [0, 1] \\
    \text{type} &= \text{softmax}\big(\text{MLP}_{d \to 32 \to 5}(\bh)\big) \in \R^5
\end{align}

Os 5 tipos de ambiguidade: PRONOUN, SCOPE, LEXICAL, STRUCTURAL, REFERENTIAL.

\paragraph{Par\^ametros totais:} ~99.4K.

\subsubsection{PlasticityGate}
\label{sec:plasticity}

Modela a plasticidade restante do modelo ao longo da sequ\^encia, simulando
deple\c{c}\~ao de recursos cognitivos:

\begin{align}
    \text{cost}_t &= \text{Softplus}\big(\text{MLP}_{d \to 16 \to 1}(\bh_t)\big) \\
    \text{budget}_t &= 1.0 - 0.3 \cdot \frac{t}{T} \quad \text{(or\c{c}amento posicional)} \\
    \text{plast}_t &= \text{clamp}\Big(\text{budget}_t - \delta \cdot \sum_{k=1}^{t} \text{cost}_k, \ 0, \ 1\Big) \\
    \text{gate}_t &= \sigma\big(\text{MLP}_{2 \to 8 \to 1}(\text{plast}_t, s_t)\big)
\end{align}

com $\delta = 0.02$ taxa de deple\c{c}\~ao.

\paragraph{Par\^ametros:} ~12.3K.

\subsubsection{MPL: Densidade e Fronteira}
\label{sec:mpl}

\paragraph{DensityTracker.}
Mant\'em um grid hash esparso 5D que estima a densidade de visita\c{c}\~ao
em cada regi\~ao do manifold. Opera com \texttt{@torch.no\_grad}:

\begin{equation}
    \rho(\bm{c}) = \frac{1}{N} \sum_{i=1}^{N} K_h(\bm{c} - \bm{c}_i)
\end{equation}

com kernel Gaussiano $K_h$ e bandwidth $h = 0.15$.

\paragraph{FrontierHead.}
Identifica regi\~oes inexploradas do manifold:

\begin{equation}
    f_{\text{frontier}} = \sigma\big(\text{MLP}_{6 \to 16 \to 1}([\bm{c}; \rho(\bm{c})])\big) \cdot (1 - \rho(\bm{c}))
\end{equation}

\paragraph{Par\^ametros:} ~129 (FrontierHead) + 0 (DensityTracker).

% -------------------------------------------------------
\subsection{Tier 3: Componentes Avan\c{c}ados (~250K params)}
% -------------------------------------------------------

\subsubsection{MOPsi: M\'odulo Orientador de $\psi$}
\label{sec:mopsi}

Modela o estado ``humano'' 5D e medeia o conflito $\phi$-$\psi$:

\begin{align}
    \bpsi &= \sigma\big(\text{MLP}_{d \to 32 \to 5}(\bh)\big) \in [0, 1]^5 \\
    \psi_{\text{score}} &= \sigma\big(\text{MLP}_{10 \to 32 \to 1}([\bpsi; \bm{c}])\big) \in [0, 1]
\end{align}

O \texttt{PhiPsiMediator} resolve conflitos:

\begin{equation}
    \text{mediation} = \sigma\big(\text{MLP}_{3 \to 8 \to 1}(\phi, \psi_{\text{score}}, \gamma)\big)
\end{equation}

\paragraph{Par\^ametros:} ~25.3K + ~41.

\subsubsection{CausalState: Condicionamento Causal}
\label{sec:causal}

Permite condicionar a gera\c{c}\~ao em um vetor de estado externo
$\bm{z} \in \R^4$:

\begin{equation}
    \bh' = \bh + g \cdot \text{MLP}_{4 \to 32 \to d}(\bm{z})
\end{equation}

onde $g = \sigma(\bm{w}^\top \bm{z} + b) \in [0, 1]$ \'e um gate
aprendido que depende apenas do vetor de estado $\bm{z}$ (n\~ao de
$\bh$), garantindo que a intensidade do condicionamento seja
determinada puramente pelo estado externo. O \texttt{PolicyBinding}
mapeia estados para par\^ametros de gera\c{c}\~ao (temperatura,
\textit{max\_tokens}) via regras heur\'isticas.

\paragraph{Par\^ametros:} ~25.5K (StateConditioning) + 0 (PolicyBinding).

\subsubsection{Metacognitive: Avalia\c{c}\~ao Contrastiva}
\label{sec:metacognitive}

O \textit{ContrastiveHead} projeta os \textit{hidden states} em duas
``vis\~oes'' e mede diverg\^encia:

\begin{align}
    \bm{v}_a &= \text{MLP}_{d \to 128 \to d/2}(\bh) \\
    \bm{v}_b &= \text{MLP}_{d \to 128 \to d/2}(\bh) \\
    \text{div} &= \sigma\big(\text{MLP}_{d \to 32 \to 1}([\bm{v}_a; \bm{v}_b])\big) \in [0, 1]
\end{align}

Alta diverg\^encia indica que o modelo tem ``duas opini\~oes'' sobre o
token, sinalizando incerteza metacognitiva.

\paragraph{Par\^ametros:} ~200K.

\subsection{Overhead Total}

\begin{table}[H]
\centering
\caption{Par\^ametros dos m\'odulos de extens\~ao.}
\label{tab:overhead}
\begin{tabular}{lrl}
\toprule
\textbf{M\'odulo} & \textbf{Params} & \textbf{\% do Modelo} \\
\midrule
EidosDecay & 368 & <0.001\% \\
Filosofia3 & 373 & <0.001\% \\
SelfModel & 982 & <0.001\% \\
Grounding (2 heads) & 99.4K & 0.028\% \\
PlasticityGate & 12.3K & 0.003\% \\
MPL (Frontier) & 129 & <0.001\% \\
MOPsi & 25.3K & 0.007\% \\
CausalState & 25.5K & 0.007\% \\
Metacognitive & 200K & 0.056\% \\
\midrule
\textbf{Total} & \textbf{~364K} & \textbf{<0.3\%} \\
\bottomrule
\end{tabular}
\end{table}

Todos os m\'odulos s\~ao opcionais e podem ser ativados/desativados via
configura\c{c}\~ao sem alterar a arquitetura base do transformer.
